\section{Retrieval and adaptive repair need separate evidence and acceptance contracts}
\label{exec19:ragrepair-scope}
A causal banking assistant needs to retrieve the right evidence, translate it
into a defensible intervention question, produce an executable analysis, and
recognize when the result does not answer that question. Retrieval and iterative
repair can support these stages, but their success criteria differ. A faithfully
quoted passage can be scientifically inapplicable; a program can pass the tests
it has repeatedly seen while using the wrong time window or treatment. The
proposal must therefore preserve the chain from source identity through causal
meaning to held-out decision quality.

Two historical originals help sharpen this design. Gao and colleagues' survey
is inspected at arXiv version 5 of 27 March 2024, distinct from its December 2023
initial discovery date. The CoT-SelfEvolve study is inspected at arXiv version 1
of 28 August 2024. Its first page and running footer also contain a 2025 conference
template and placeholder DOI; those do not establish a verified proceedings
version. Archive metadata lists its first author as Quoc, Thai Tang, while the
PDF displays Thai Tang Quoc. The version and author-name distinction is retained
in the audit. Neither original establishes complete October 2026 frontier coverage.
\citep{exec19_rag_survey,exec19_selfcorrection_study}

\subsection{Retrieval faithfulness is narrower than causal evidence}
\label{exec19:ragrepair-retrieval}
The survey distinguishes a basic index--retrieve--generate pipeline, advanced
pre- and post-retrieval improvements, and modular pipelines that can rearrange
or repeat retrieval and generation. Basic indexing converts source documents
into chunks and representations; retrieval ranks context for a query; generation
uses that context to answer. Query rewriting, finer indexing, metadata, reranking
and compression can improve the fit between question and material. Modular
search and iterative retrieval provide alternatives when a single initial
query is inadequate. These are useful architectural choices, not proof that a
retrieved statement is true or that a generated graph is identified.
\citep{exec19_rag_survey}

For a bank asking whether a servicing call prevents default, retrieval must
retain more than a passage containing both terms. An observational prediction
study, a randomized call policy and a vendor's marketing description can all
be semantically similar to the question while supporting different claims.
The proposed evidence object therefore records source type, original date and
version, population, treatment and outcome definitions, assignment mechanism,
measurement clock, assumptions and the precise claim supported. Rewriting the
business question may produce several candidate estimands; it must not silently
replace an intervention with a correlation because that version retrieves more
text. This evidence-object design is a project proposal motivated by the survey's
retrieval and augmentation distinctions, not a capability reported for an
already integrated package.
\citep{exec19_rag_survey}

The survey separately evaluates retrieval quality and generation quality.
It distinguishes context relevance, answer faithfulness and answer relevance,
and discusses noise resistance, refusal when necessary knowledge is missing,
integration across documents and resistance to known false retrieved statements.
Its term counterfactual robustness means recognizing and rejecting inaccuracies
in documents. It does not mean identifying a potential outcome or preserving
cross-world structural dependence. Confusing these uses of the same word would
let a language benchmark appear to qualify a causal capability it never tested.
The survey also warns that the listed evaluation metrics were not yet a mature,
standardized system. A package name or metric family in its comparison is a
lead for primary inspection, not an automatically adopted scoring implementation.
\citep{exec19_rag_survey}

A public evaluation can make these distinctions concrete before a costly bank
trial. Construct matched evidence packets containing a relevant original,
a wrong-population study, a later revision, a contradictory source and an
unsupported confident claim. Preserve their identities and the gold-supported
conclusions. Score retrieval of the needed evidence separately from faithful
quotation, applicability to the declared estimand and justified refusal.
Then require the proposed graph to mark unsupported edges and alternatives.
Removing an irrelevant source should not alter a justified answer; removing
the only identification assumption should change its status. These are proposed
perturbation and evidence tests, not a claim that a retrieved paper establishes
an actual bank's assignment mechanism. The survey's discussion of noise,
contradiction and source/metadata exposure motivates such explicit checks.
\citep{exec19_rag_survey}

\subsection{What the self-correction study actually measures}
\label{exec19:ragrepair-study}
CoT-SelfEvolve combines retrieved StackOverflow discussions, generated reasoning
prompts, initial code, syntax checking, execution against unit tests and feedback-
guided revisions. Its feedback is partly external: tracebacks and test outcomes
come from an executor rather than from a model merely reconsidering its prose.
The inspected protocol uses the Completion setting of DS-1000, comprising
1,000 problems across seven Python libraries. It reports a retrieval collection
of 558,402 posts and 972,513 comments, with a minimum-comment filter and a token
budget for assembled documents. The paper also describes its benchmark tasks
as extracted from StackOverflow. This shared origin motivates an explicit
retrieval/benchmark overlap audit; the inspected descriptions do not establish
that leakage occurred or that every relevant overlap was removed.
\citep{exec19_selfcorrection_study}

The study defines pass@$n$ as the fraction of problems for which all tests pass
within at most $n$ sequential attempts, examining budgets from one to five.
Later attempts can use feedback from earlier ones. This differs from the event
that a uniformly chosen subset of an independent candidate pool contains a
passing program. It also differs from passing a fresh causal acceptance test.
The proposed product should retain the paper's operational definition rather
than treating identical-looking pass notation as a common metric.
\citep{exec19_selfcorrection_study}

The reported comparisons are informative but protocol-specific. Table 1 gives
an aggregate change from 30.5 to 35.5 percent when both prompt generators are
enabled; Table 2 gives 35.5 to 39.5 percent when the prompt generator uses the
larger model while the code generator remains fixed. Arithmetic on those
printed, rounded values gives
\begin{align}
 35.5-30.5&=5.0\text{ percentage points},&
 \frac{35.5-30.5}{30.5}&=\frac{10}{61},\nonumber\\
 39.5-35.5&=4.0\text{ percentage points},&
 \frac{39.5-35.5}{35.5}&=\frac{8}{71}.
 \label{exec19:ragrepair-percentages}
\end{align}
The relative changes are approximately 16.39 and 11.27 percent, respectively;
they are not increases of that many percentage points. Table 1 also shows
heterogeneity: the SciPy entry changes from 33.02 to 32.08 percent with both
prompt generators enabled. A positive aggregate therefore does not establish
a gain for every library or every causal-analysis task. No raw counts or paired
uncertainty interval is recovered from these rounded percentages.
\citep{exec19_selfcorrection_study,exec19_retrieval_repair_reference}

For its GPT-4 attempt-budget experiment, Table 3 reports overall success rising
from 14.0 to 83.2 percent between one and five allowed attempts. The text reports
prompt tokens rising from 2.54 million to 9.29 million, while completion tokens
do not rise at the same rate. This is evidence about the reported historical
protocol, not a present model benchmark or a full cost-benefit analysis.
Cross-figure and cross-table totals should not be pooled without reconciling
settings and runs. The authors themselves discuss proprietary model drift,
temperature and sampling settings, and the dependence on available unit tests.
Their analogy between five refinement steps and five usability testers does
not supply an optimal-stopping theorem for code or causal analysis.
\citep{exec19_selfcorrection_study}

Executable reuse has an additional unresolved prerequisite. The selected
original's reproducibility statement gives \texttt{ANONYMIZED-LINK} for its code
and prepared data. That is not a resolved repository identity, version or
license. It is therefore appropriate to retain the method as a research
comparator while deferring code adoption. A similarly named repository cannot
be substituted without establishing authorship and correspondence. Proposed
CI integration in the discussion also does not demonstrate that such an
integration has already been built here.
\citep{exec19_selfcorrection_study}

\subsection{Count failed attempts and accepted causal outputs}
\label{exec19:ragrepair-accounting}
A small authored example shows why the stopping and acceptance records matter.
Consider ten requests with a maximum of three attempts. Their first test-passing
attempts are
\begin{equation}
 (1,1,2,2,2,3,\varnothing,\varnothing,\varnothing,\varnothing),
 \label{exec19:ragrepair-attempts}
\end{equation}
where $\varnothing$ means that no candidate passes within the budget. The
cumulative counts of passing requests are two, five and six. The six successful
requests use eleven attempts in total; the four failures consume twelve more.
Thus all-request cost uses 23 attempts, not eleven. The average attempt count
is 2.3 across all requests, compared with $11/6$ if failures are omitted.
Suppose a separate causal review accepts only four of the six test-passing
outputs because two violate the frozen estimand or information clock. Test
success is $6/10$, while accepted causal output is $4/10$.
\citep{exec19_retrieval_repair_reference}

Assume USD 2 per attempt and USD 15 value per accepted output. The complete
cost and net value are
\begin{align}
 C&=2(23)=\mathrm{USD}\ 46,\nonumber\\
 C/4&=\mathrm{USD}\ 11.50\text{ per accepted output},\nonumber\\
 V_{\mathrm{accepted}}&=15(4)-46=\mathrm{USD}\ 14,\nonumber\\
 V_{\mathrm{test\ only}}&=15(6)-46=\mathrm{USD}\ 44.
 \label{exec19:ragrepair-cost}
\end{align}
Equating test passes with accepted causal outputs overstates this stipulated
net value by USD 30. The example does not measure actual token prices,
engineering labor, model performance or bank ROI. Its role is to enforce a
complete denominator before real economic evaluation, including unsuccessful
retrieval, repair, runtime and human adjudication.
\citep{exec19_retrieval_repair_reference}

\subsection{A public prototype protocol before paid trials}
\label{exec19:ragrepair-protocol}
The proposal can now specify a bounded public experiment. Freeze the source
corpus and exclusions, model/backend version, prompt and package environment,
maximum attempts and stopping rule. Compare one-shot generation, executor-
feedback repair and retrieval-assisted repair under the same request budget.
Keep the feedback tests distinct from a final acceptance set that checks causal
target, time alignment, graph and identification scope, uncertainty and refusal.
Include deliberate missing-evidence cases, misleading retrieval, future leakage,
wrong treatment definitions and syntactically correct but causally invalid
programs. Log every retrieved source and every attempt; assess accepted outputs
and complete cost jointly. This design combines the survey's component-level
criteria with the empirical study's adaptive feedback structure without
borrowing either source's reported score as a product guarantee.
\citep{exec19_rag_survey,exec19_selfcorrection_study,exec19_retrieval_repair_reference}

Proceeding requires a genuine artifact for an executable CoT-SelfEvolve
comparison, inspected originals for any adopted retrieval/evaluation package,
and an independently specified final acceptance oracle. Where these are absent,
use explicitly authored reference components and label them accordingly.
Public evidence can justify whether a prototype merits an expensive banking
experiment; it cannot establish institutional data truth or replace subsequent
owner acceptance. The historical findings narrow the design and expose concrete
failure modes while leaving broader frontier coverage and adoption open.
\citep{exec19_rag_survey,exec19_selfcorrection_study}
